\section{Ambient dimensions and discrete consequences}\label{sec:dimension}
The proofs in Appendices~\ref{app:positive} and~\ref{app:negative} take place
in $d=m-1$. We use the following two external results with different feasible
classes.

Heilman--Tarter \cite[Theorem 1.9]{HT2021} states that, for positive
correlation and prescribed positive cell masses, every maximizing Gaussian
partition in dimension $N\geq m-1$ depends, after a common orthogonal change
of coordinates, on at most $m-1$ coordinates.
Heilman \cite[Problem 7.1 and Theorem 7.9]{H2022} instead treats minimizers of
a positive-kernel bilinear objective over two partitions whose corresponding
masses agree but are \emph{free}. Every such minimizer has a common
$(m-1)$-dimensional reduction. This latter result has no hyperstability
hypothesis and is not the distinct theorem for two separately prescribed
mass vectors.

\subsection{Existence and the positive problem}
For finite $N$, the simplex-valued feasible set with prescribed masses is
weakly compact in $L^2(\gamma_N)^m$. The Ornstein--Uhlenbeck operator is compact:
each Hermite degree has finite multiplicity and $|\rho|^k\to0$.
Thus the positive quadratic objective is weakly continuous on this bounded
set and attains its maximum.

Every maximizer is an indicator partition. Otherwise two fixed labels are
both uniformly fractional on a positive-measure set. Choose two disjoint
subsets of equal positive measure there. Opposite label perturbations on
those subsets preserve all masses and produce two distinct feasible
assignments whose average is the original one. Strict positivity of $T_\rho$
for $\rho>0$, visible from its strictly positive Hermite eigenvalues, makes
the average of their objectives strictly larger. This contradicts maximality.
Theorem~1.9 of \cite{HT2021} therefore applies to every relaxed maximizer.
Appendix~\ref{app:positive} gives its value and its exact equality form.

\subsection{Free matching masses and the negative problem}
The set of pairs with free common masses is weakly compact, and compactness
of $T_{-r}$ makes the bilinear objective jointly weakly continuous on it.
A minimizing pair exists. Hold its attained mass vector $p$ fixed.
With $h$ fixed, the minimizing face for the continuous linear functional in
$f$ has an extreme point. Every extreme point of this fixed-mass
simplex-assignment set is an indicator, by the same equal-measure opposite
perturbation construction. Choose such an $f_1$, and then with $f_1$ fixed
choose an indicator minimizer $h_1$. The pair has the same global minimum.
This sequential argument does not require an incorrect general principle
about extreme points of bilinear optimization.

The competitive bounds~\eqref{eq:competitive-negative} hold in every finite
ambient dimension, and~\eqref{eq:negative-range} makes all masses positive.
Reflect the second partition: $B_{-r}(f_1,h_1)=B_r(f_1,h_1(-\,\cdot))$.
The free-common-mass theorem \cite[Theorem 7.9]{H2022} now reduces this
indicator minimizer to a common subspace of dimension at most $d$.
Reflection back preserves that subspace. Appendix~\ref{app:negative} forces
$f_1=h_1=g$ in one common frame.

For the original fractional pair, the intermediate pair $(g,h)$ still has
the minimum value. Its excess over $(g,g)$ is exactly $D_-(h)$.
Strict reference ordering forces $h=g$. The original pair $(f,g)$ then
has excess $D_-(f)$, so $f=g$. This recovers equality for the original
fractional pair, not only for an indicator representative.
In either sign, an assignment in $N<d$ can be embedded in dimension $d$;
equality would require a regular centroid Gram of rank $d$, impossible for
the embedded rank at most $N$. This completes Theorem~\ref{thm:noise}.

\subsection{The precise low-influence corollary}
On $[m]^n$ use uniform product measure and the noise channel which, in each
coordinate, keeps a symbol with probability
$(1+(m-1)\rho)/m$ and changes it to each particular different symbol with
probability $(1-\rho)/m$. It exists for $-1/(m-1)\leq\rho\leq1$.
For $f:[m]^n\to\Delta_m$ define
\[
 \operatorname{Inf}_k(f)=\sum_i\E\Var(f_i\mid\hbox{all coordinates except }k).
\]
Write $\operatorname{Stab}_\rho(f)=\sum_i\E[f_i(X)f_i(Y)]$ for this channel,
and $p_m(\rho)=S_\rho(g)$.

\begin{corollary}\label{cor:discrete}
Fix $m\geq3$, an admissible nonzero $\rho$ in the corresponding interval of
Theorem~\ref{thm:noise}, and $\eps>0$. There is $\tau>0$, depending on
$m,\rho,\eps$, such that for every $n$ and every $f$ with
$\max_k\operatorname{Inf}_k(f)\leq\tau$,
\[
 \operatorname{Stab}_\rho(f)\leq p_m(\rho)+\eps
 \quad(\rho>0,\ \E f=\1/m),
\]
whereas
\[
 \operatorname{Stab}_\rho(f)\geq p_m(\rho)-\eps
 \quad(\rho<0,\ \text{with arbitrary }\E f).
\]
\end{corollary}
\begin{proof}
Isaksson--Mossel \cite[Theorem 7.1]{IM2009} transfers a low-influence
simplex-valued function to a Gaussian relaxed assignment with the same
expectation vector and approximates stability at the same fixed $\rho$.
Apply Theorem~\ref{thm:noise}, using its free-mass single-assignment consequence
for the negative case, and allocate the approximation error below $\eps$.
Lemma~7.3 of \cite{IM2009} identifies $p_m(\rho)$ with the limiting plurality
stability. This gives the stated additive comparison.
\end{proof}
There is no uniform-in-$\rho$ influence threshold or assertion of exact
finite-electorate optimality. Only the single-rule arbitrary-bias negative
corollary is asserted here; no additional discrete bilinear theorem is needed.
The near-zero Gaussian intervals do not give a full-range Standard Simplex
theorem or a sharp MAX-$m$-CUT hardness ratio.

